\SourceLeaf{401}%
\Needspace{10\baselineskip}
\begin{center}\bfseries [6.]\end{center}
\noindent\hspace*{1.2em}Ist $df$ das Differential der Seitenkrümmung, welches dem Differential $du$ entspricht, so ist in den Zeichen der Entwickelung:
\[
\begin{aligned}
&2(Pt't'+2Qt'\theta'+R\theta'\theta')\BareRoot{(PR-QQ)}.df\\
&\quad=2(PR-QQ)(t'd\theta'-\theta'dt')\\
&\qquad+t'^3(-QP'+2PQ'-PP'')du\\
&\qquad+t't'\theta'(-RP'+2QQ'+2PR'-3QP'')du\\
&\qquad+t'\theta'\theta'(3QR'-2RP''-2QQ''+PR'')du\\
&\qquad+\theta'^3(RR'-2RQ''+QR'')du.
\end{aligned}
\]
Wenn man bloss das erste Glied des zweiten Theils der Gleichung berücksichtigt, ist das Integral
\[
f=\operatorname{Arc\,tang}\frac{\theta'\BareRoot{(PR-QQ)}}{Pt'+Q\theta'}+\text{const.}=\chi+\text{const.}
\]
Aus der vollständigen Differentiation von $\chi$ erhält man aber (weil
\[
2\BareRoot{b}(cc+aab).d\chi=2bc.da+ac.db-2ab.dc,
\]
\[
\operatorname{tang}\chi=\frac{a\BareRoot{b}}{c}\quad\text{gesetzt,)}
\]
[die Gleichung:
\[
\begin{aligned}
&2P(Pt't'+2Qt'\theta'+R\theta'\theta')\BareRoot{(PR-QQ)}d\chi\\
&=2(PR-QQ)(Pt'+Q\theta')d\theta'+\theta'(Pt'+Q\theta')d(PR-QQ)\\
&\quad-2\theta'(PR-QQ)d(Pt'+Q\theta')
\end{aligned}
\]
oder]
\[
\begin{aligned}
&2(Pt't'+2Qt'\theta'+R\theta'\theta')\BareRoot{(PR-QQ)}d\chi\\
&\quad=2(PR-QQ)(t'd\theta'-\theta'dt')\\
&\qquad+((2QQ-PR)t'\theta'+QR\theta'\theta').\frac{dP}{P}\\
&\qquad-2(Qt'\theta'+R\theta'\theta').dQ\\
&\qquad+(Pt'\theta'+Q\theta'\theta').dR.
\end{aligned}
\]
